\subsection{Geometria, discretização e volume}\label{sec:geometria}

A formulação de dimensionamento reaproveita o núcleo mecânico do manuscrito da plataforma, baseado no roteiro de pontes de madeira de \textcite{CALIL2006}, e explicita as hipóteses necessárias à comparação entre materiais. As longarinas são representadas por vigas prismáticas circulares; as pranchas do tabuleiro vencem transversalmente os vãos entre longarinas. Com $B$ denotando a largura da pista, as propriedades das seções são
\begin{equation}
 A_\ell=\frac{\pi d^2}{4},\quad I_\ell=\frac{\pi d^4}{64},\quad
 W_\ell=\frac{\pi d^3}{32},\qquad
 A_t=b_wh,\quad I_t=\frac{b_wh^3}{12},\quad W_t=\frac{b_wh^2}{6}.
\label{eq:secoes}
\end{equation}
O diâmetro $d$ é o diâmetro uniforme de cálculo. A representação de uma tora afilada por um diâmetro equivalente deverá ser compatibilizada com o critério normativo aplicável; o volume de uma peça prismática equivalente não é necessariamente o volume comercial da tora afilada.

A disposição construtiva considera números inteiros de peças. Para um comprimento disponível $D$, largura de peça $w$ e espaçamento livre proposto $s$, a rotina geométrica examina os inteiros adjacentes a $(D+s)/(w+s)$ e calcula
\begin{equation}
 s_{corr}=\frac{D-nw}{n-1},\qquad n\geq2.
\label{eq:acomodacao}
\end{equation}
Empregam-se $(D,w)=(B,d)$ para as longarinas e $(D,w)=(L,b_w)$ para as pranchas. Entre as alternativas, priorizam-se a viabilidade dos espaçamentos e a proximidade ao valor proposto. Essa operação introduz descontinuidades no espaço de projeto e permite que o número de peças varie durante a busca. O problema, portanto, não é uma simples otimização contínua de seções com número de elementos fixo.

O volume de madeira contabilizado é
\begin{equation}
 V(\mathbf{x})=n_\ell A_\ell L+n_t A_t B.
\label{eq:volume_modelo}
\end{equation}
Esse indicador cobre longarinas e pranchas, sem representar diretamente custo, carbono incorporado, perdas de corte ou volume de fundações e ligações.

\subsection{Ações e distribuição transversal}\label{sec:acoes}

A configuração de referência transportada do manuscrito da Matéria utiliza o veículo TB-240 associado à \textcite{nbr7188}: seis rodas de $P=40$~kN, em três eixos espaçados de $a=1{,}5$~m, correspondendo a 80~kN por eixo e 240~kN por veículo. A carga de multidão de referência é $q_A=4$~kN/m\textsuperscript{2}. A aplicação dessa configuração à via estudada e as combinações correspondentes deverão ser documentadas na definição final dos casos.

Para peso específico $\gamma_{mat}=9{,}81\rho_{ap}/1000$ em kN/m\textsuperscript{3}, o peso próprio médio do tabuleiro por unidade de área é $g_{A,t}=\gamma_{mat,t}n_tA_t/L$. Definindo $b_{tr,i}$ como largura tributária da longarina $i$ e $g_{A,ad}$ como carga permanente adicional, tem-se
\begin{equation}
 g_{\ell,i}=\gamma_{mat,\ell}A_\ell+(g_{A,t}+g_{A,ad})b_{tr,i},
 \qquad q_{\ell,i}=q_A b_{tr,i}.
\label{eq:acoes_lineares}
\end{equation}
A largura tributária tem dimensão de comprimento e deve ser definida a partir da distribuição efetiva das cargas e das posições dos apoios. O espaçamento livre $s_{corr}$, a distância entre eixos das longarinas $s_{corr}+d$ e o vão de cálculo do tabuleiro $\ell_t$ são grandezas distintas. Para cargas concentradas, a parcela recebida pela longarina depende também da posição transversal da roda; a adoção de uma roda inteira sobre uma longarina constitui uma hipótese de carregamento a justificar, e não uma consequência da conversão de carga de área em carga linear.

\begin{fillblock}
\noindent\textbf{[COMPATIBILIZAÇÃO IDENTIFICADA NA TRANSFERÊNCIA]} A implementação atual utiliza o espaçamento livre corrigido como largura tributária e como vão do tabuleiro. Antes das novas rodadas, conferir a representação dos apoios, as longarinas de borda e o balanço de cargas, exigindo que a distribuição preserve a carga total. A formulação acima distingue essas grandezas; sua adoção definitiva requer compatibilização com o código. Esta revisão do manuscrito não altera o programa de cálculo.
\end{fillblock}

\subsection{Esforços e resposta elástica}\label{sec:esforcos}

Adotam-se pequenas deformações e comportamento elástico linear de vigas de Euler--Bernoulli, sem deformação por cisalhamento. Para carga permanente uniforme $g_\ell$, as expressões são
\begin{equation}
 M_{G,k}=\frac{g_\ell L^2}{8},\qquad
 Q_{G,k}=\frac{g_\ell L}{2},\qquad
 \delta_{G,k}=\frac{5g_\ell L^4}{384E_{M0}I_\ell}.
\label{eq:permanentes_modelo}
\end{equation}
No cálculo em kN e m, os módulos e resistências são convertidos de MPa para kN/m\textsuperscript{2}, multiplicando-se por $10^3$.

Para a posição simétrica de três cargas iguais, com uma carga no meio do vão e as demais afastadas de $a$, o momento e o deslocamento no meio do vão são
\begin{align}
 M_{P,k}(L/2)&=\frac{3PL}{4}-Pa,\label{eq:momento_tres_rodas}\\
 \delta_{P,k}(L/2)&=\frac{P}{48E_{M0}I_\ell}
 \left[L^3+2b(3L^2-4b^2)\right],\qquad b=\frac{L-2a}{2}.
\label{eq:flecha_tres_rodas}
\end{align}
Essas expressões pressupõem que as três cargas estejam sobre o vão, com $L\geq2a$; $a$ é o espaçamento entre eixos do veículo. São soluções de uma configuração especificada e não substituem, sem verificação, a envoltória de todas as posições possíveis do veículo.

A conferência independente da carga móvel será feita por linhas de influência. Para uma seção $x$ e uma carga na posição $z$, a linha de influência do momento da viga biapoiada é
\begin{equation}
 \eta_M(x,z)=
 \begin{cases}
 z(L-x)/L,&0\leq z\leq x,\\
 x(L-z)/L,&x<z\leq L.
 \end{cases}
\label{eq:linha_influencia}
\end{equation}
Para cada posição $\zeta$ do veículo, somam-se as rodas que efetivamente estão sobre o vão e integra-se a carga distribuída sobre a região carregada $\Omega(\zeta)$:
\begin{equation}
 M_Q(x;\zeta)=\sum_{j:z_j(\zeta)\in[0,L]}P_j\eta_M(x,z_j)
 +\int_{\Omega(\zeta)}q_\ell(z)\eta_M(x,z)\,dz.
\label{eq:carga_movel}
\end{equation}
A envoltória resulta da busca em $x$ e $\zeta$; cortantes e deslocamentos são conferidos pelos respectivos operadores de resposta. A posição mais desfavorável pode ser diferente para cada grandeza.

O coeficiente de impacto vertical é aplicado uma única vez aos efeitos variáveis para os quais for prescrito. A rotina existente incorpora o impacto nos esforços resistentes, mas calcula a flecha móvel pelas cargas estáticas das rodas, sem parcela de multidão. A composição da ação variável no ELS, incluindo multidão e eventual amplificação dinâmica, deverá ser verificada antes de interpretar diferenças de desempenho. Para a comparação de materiais, a regra escolhida será idêntica nos dois cenários.

\begin{fillblock}
\noindent\textbf{[DOMÍNIO DO MODELO HERDADO]} Conferir as expressões simplificadas por trechos com a envoltória da Equação~\ref{eq:carga_movel}. A rotina de cortante usa uma grandeza auxiliar $e=L-3a-2d$, que pode ser negativa nos vãos curtos da campanha; não interpretar um comprimento negativo como trecho carregado. Conferir também o tratamento de $L=10$~m no coeficiente de impacto, pois a fronteira dos trechos difere entre texto e código. Essas verificações precedem a execução da nova matriz de casos.
\end{fillblock}

\subsection{Resistência, fluência e restrições normalizadas}\label{sec:restricoes}

As resistências de cálculo são expressas por
\begin{equation}
 f_{m,d}=\frac{k_{mod}f_{m,k}}{\gamma_{wf}},\qquad
 f_{v0,d}=\frac{k_{mod}f_{v0,k}}{\gamma_{wc}},\qquad
 k_{mod}=k_{mod,1}k_{mod,2}.
\label{eq:resistencias_modelo}
\end{equation}
Os fatores de modificação são associados à duração do carregamento e à condição de umidade \parencite{NBR7190-1:2022}. O módulo médio empregado nas flechas não é minorado pelos mesmos fatores de resistência. A fluência é tratada separadamente.

Na combinação com uma ação variável predominante, os efeitos de cálculo são
\begin{equation}
 M_d=\gamma_gM_{G,k}+\gamma_qM_{Q,k}^{dyn},\qquad
 Q_d=\gamma_gQ_{G,k}+\gamma_qQ_{Q,k}^{dyn},
\label{eq:combinacao_modelo}
\end{equation}
com as combinações e fatores definidos para as situações de projeto consideradas \parencite{NBR8681:2025}. O sobrescrito $dyn$ explicita a inclusão do impacto, evitando sua dupla aplicação.

A formulação herdada avalia a flecha final da combinação quase permanente e a flecha variável separadamente:
\begin{equation}
 \delta_{fin,qp}=(1+\varphi)(\delta_{G,k}+\psi_2\delta_{Q,k}),\qquad
 U_{fin}=\frac{\delta_{fin,qp}}{L/250},\qquad
 U_Q=\frac{\delta_{Q,k}}{L/360}.
\label{eq:servico_modelo}
\end{equation}
A primeira parcela não deve ser denominada deslocamento máximo total sob passagem do veículo: a parcela variável é reduzida por $\psi_2$. Os limites $L/250$ e $L/360$ são os critérios transportados do modelo da Matéria e serão associados explicitamente à combinação e à disposição normativa aplicáveis. A reavaliação C$\rightarrow$R deve verificar ambos, pois o projeto pode atender a um e exceder o outro.

As quatro restrições mecânicas implementadas assumem a forma
\begin{align}
 g_1&=\frac{|M_{d,\ell}|}{W_\ell f_{m,d,\ell}}-1,&
 g_2&=\frac{4|Q_{d,\ell}|}{3A_\ell f_{v0,d,\ell}}-1,\label{eq:g12}\\
 g_3&=\max(U_{fin},U_Q)-1,&
 g_4&=\frac{|M_{d,t}|}{W_t f_{m,d,t}}-1.
\label{eq:g34}
\end{align}
A tensão máxima de cisalhamento da seção circular maciça é $4Q/(3A)$, equivalente a $QS/(Id)$ com $S=d^3/12$. Para o tabuleiro, o momento permanente da idealização biapoiada é $g_t\ell_t^2/8$. A parcela de roda depende da extensão e posição do contato sobre o vão; a expressão herdada $P(\ell_t-a_r)/4$, com $a_r=0{,}45$~m, exige conferência dessa idealização e de seu domínio antes de ser aplicada a vãos comparáveis ou inferiores ao contato.

As duas restrições geométricas $g_5$ e $g_6$ correspondem, respectivamente, aos limites de espaçamento livre das longarinas e do tabuleiro. Para cada família,
\begin{equation}
 g_{esp}=\max\left\{\frac{s_{min}-s_{corr}}{s_{min}},
                         \frac{s_{corr}-s_{max}}{s_{max}}\right\}\leq0,
\label{eq:restricao_espaco}
\end{equation}
omitindo-se o primeiro termo quando $s_{min}=0$. Configurações sem acomodação geométrica admissível são rejeitadas. Esses limites são parâmetros construtivos do estudo; não são chamados normativos sem referência específica.

As seis restrições não constituem verificação completa de uma ponte: estabilidade, apoio perpendicular às fibras, ligações, durabilidade e outros modos não representados devem ser avaliados separadamente em um projeto executivo.
